% TOP_PROBLEM: {"id": 2302013, "title": "Research Problems in Function Theory — Problem 2.13", "category_id": 8, "category": {"id": 8, "name": "analysis", "display_name": "Analysis", "description": "Limits, continuity, calculus, and function theory.", "slug": "analysis", "order_index": 8, "created_at": "2026-07-31T15:26:25.671Z"}, "status": "open"}
% FIRST_POSTED: null
% ENTRY_KIND: statement_only
% Imported from: batch11.tex; UnsolvedMath ID 2302013
% Compile twice: lualatex 2302013.tex
\documentclass[11pt]{article}
\usepackage{amsmath,amssymb,mathtools}
\usepackage{fontspec}
\usepackage{unicode-math}
\setmainfont{Latin Modern Roman}
\setmathfont{Latin Modern Math}
\usepackage{xurl}
\usepackage[unicode,colorlinks=true,linkcolor=blue,urlcolor=blue,pdftitle={UnsolvedMath Problem 2302013}]{hyperref}
\newcommand{\sref}[2]{\hyperlink{source:#1:#2}{[S#2]}}
\newcommand{\cataloguescope}[1]{\paragraph{Mathematical statement.}#1}
\begin{document}
% UnsolvedMath ID 2302013; source catalogue code AMR-022-2013
\setcounter{section}{2302012}
\section{Research Problems in Function Theory — Problem 2.13}\label{2302013}
{\small\sffamily UnsolvedMath ID 2302013 · Analysis}
\subsection{Definitions and mathematical statement}

\paragraph{Shared notation.}$\mathbb N=\{1,2,\ldots\}$, and $\mathbb Z,\mathbb Q,\mathbb R,\mathbb C$ denote the integers, rationals, reals and complex numbers. Any other conventions are specified locally.


For an entire function $f$, write
\[
 M(r,f)=\max_{|z|=r}|f(z)|,\qquad m_0(r,f)=\min_{|z|=r}|f(z)|.
\]
For nonconstant $f$, its order is $\rho(f)=\limsup_{r\to\infty}\log\log M(r,f)/\log r$;
the lower order is defined using $\liminf$ instead. Every constant entire function, including the zero function, has order and lower order zero. In logarithmic modulus ratios, use the extended-real convention $\log0=-\infty$.
Let $f(z)=a_0+\sum_{n\ge1}a_nz^{\lambda_n}$ be transcendental entire,
where $a_n\ne0$ and $0<\lambda_1<\lambda_2<\cdots$ are integer exponents.
For $a\in\mathbb C$, let $n(r,a,f)$ count zeros of $f-a$ in $|z|\le r$, with multiplicities, and put
\[
 N(r,a,f)=\int_0^r\frac{n(t,a,f)-n(0,a,f)}{t}\,dt+n(0,a,f)\log r,
 \quad T(r,f)=\frac1{2\pi}\int_0^{2\pi}\log^+|f(re^{i\theta})|\,d\theta.
\]
A finite Picard exceptional value is a value taken only finitely often. A finite Borel exceptional value
has exponent of convergence of its $a$-points smaller than $\rho(f)$: that exponent is
$\inf\{s>0:\sum_{f(z)=a,\ z\ne0}|z|^{-s}<\infty\}$, counting multiplicity and taking the infimum to be $+\infty$ if appropriate.
A finite deficient value has $\delta(a,f)=1-\limsup_{r\to\infty}N(r,a,f)/T(r,f)>0$.

\paragraph{Statement recorded in the supplied source.}
If $\lambda_n/n\to\infty$, must $f$ have (a) no finite Picard exceptional value,
(b) no finite Borel exceptional value, and (c) no finite deficient value?
If (b) or (c) fails under this gap condition, does that assertion hold under the stronger condition
$\sum_{n\ge1}\lambda_n^{-1}<\infty$? \sref{2302013}{2}
The 2018 update records Murai's counterexample to the weaker gap condition for deficiency and his positive results for the reciprocal-summability condition in (b) and (c). The original problem also attributes the Picard conclusion under reciprocal summability to Biernacki. \sref{2302013}{2}
\subsection{Short English statement}
Determine which kinds of exceptional finite values are excluded by large gaps in an entire function's Taylor series, and whether summable reciprocal exponents give stronger conclusions.
\subsection{Sources}
\begin{description}
\item[\hypertarget{source:2302013:1}{S1}] ULAM.AI and the UnsolvedMath Contributors, \emph{UnsolvedMath: A Curated Collection of Open Mathematics Problems}, record 2302013 (AMR-022-2013). CC BY 4.0. \url{https://huggingface.co/datasets/ulamai/UnsolvedMath}.
\item[\hypertarget{source:2302013:2}{S2}] W. K. Hayman and E. F. Lingham, \emph{Research Problems in Function Theory} (2018), Chapter 2, Problem 2.13 and Update 2.13; Chapter 1 notation. \url{https://arxiv.org/abs/1809.07200}.
\end{description}
\label{end:2302013}
\end{document}
